\documentclass[11pt]{letter}
\usepackage[utf8]{inputenc}
\usepackage[T1]{fontenc}
\usepackage{mathptmx}
\usepackage[margin=2.5cm]{geometry}
\usepackage{xcolor}
\usepackage{hyperref}
\usepackage{setspace}
\onehalfspacing

\hypersetup{
  colorlinks=true,
  linkcolor=blue!70!black,
  urlcolor=blue!70!black,
}

\signature{Lucas Rover\\
Programa de P\'os-Gradua\c{c}\~ao em Engenharia Mec\^anica\\
Universidade Tecnol\'ogica Federal do Paran\'a (UTFPR)\\
Ponta Grossa, Paran\'a, Brasil\\
\href{mailto:lucasrover@utfpr.edu.br}{lucasrover@utfpr.edu.br}}

\address{Lucas Rover\\
Universidade Tecnol\'ogica Federal do Paran\'a (UTFPR)\\
Ponta Grossa, Paran\'a, Brasil}

\begin{document}

\begin{letter}{Aos Editores\\
\textit{Nature Machine Intelligence}}

\opening{Prezados Editores,}

Submetemos \textbf{``Same Prompt, Different Answer: Exposing the Reproducibility Illusion in Large Language Model APIs''} para considera\c{c}\~ao como Artigo na \textit{Nature Machine Intelligence}.

Mais de 1,5 bilh\~ao de pessoas interagem semanalmente com grandes modelos de linguagem por meio de APIs comerciais. Somente o ChatGPT atende mais de 800 milh\~oes de usu\'arios semanais, o Gemini ultrapassa 750 milh\~oes mensais e o Claude alcan\c{c}a 19 milh\~oes com crescimento anual de 190\%. Todo grande provedor de API oferece um par\^ametro de temperatura que, quando definido como zero, \'e documentado como produzindo sa\'idas determin\'isticas. Pesquisadores de medicina, engenharia, ci\^encias ambientais e an\'alise de dados dependem dessa garantia sempre que reportam resultados gerados por LLMs. Nosso manuscrito mostra que essa garantia n\~ao se sustenta.

Conduzimos 4.104 experimentos controlados com oito modelos e cinco provedores de API. Sob decodifica\c{c}\~ao gulosa com temperatura zero e sementes aleat\'orias fixas, modelos implantados localmente reproduziram suas pr\'oprias sa\'idas 96\% das vezes, enquanto modelos servidos por API alcan\c{c}aram apenas 22\%. A diferen\c{c}a supera quatro vezes, e \'e invis\'ivel aos usu\'arios: as sa\'idas permanecem semanticamente equivalentes (BERTScore F1~$>$~0,97), mas diferem textualmente, de modo que um pesquisador lendo duas respostas lado a lado n\~ao perceberia a varia\c{c}\~ao. Um experimento de quase-isolamento, no qual os mesmos pesos do modelo foram servidos por um provedor em nuvem com infraestrutura m\'inima, recuperou reprodutibilidade pr\'oxima \`a local, atribuindo a varia\c{c}\~ao \`a complexidade da infraestrutura de produ\c{c}\~ao (paralelismo de tensores, batching din\^amico, decodifica\c{c}\~ao especulativa) e n\~ao \`a implanta\c{c}\~ao em nuvem em si. Complementamos esse diagn\'ostico com um protocolo leve de proven\^ancia que adiciona menos de 1\% de sobrecarga computacional e torna a varia\c{c}\~ao oculta detect\'avel.

Trabalhos anteriores abordaram esse problema apenas parcialmente. Atil et al.\ (2024) testaram o ChatGPT isoladamente; Yuan et al.\ (NeurIPS 2025) focaram na precis\~ao num\'erica de um \'unico provedor. Nosso estudo \'e o primeiro a abranger cinco provedores e tr\^es implanta\c{c}\~oes locais, estender a an\'alise a fluxos multi-turno e de gera\c{c}\~ao aumentada por recupera\c{c}\~ao (RAG), e entregar tanto atribui\c{c}\~ao causal quanto uma ferramenta de mitiga\c{c}\~ao pronta para uso, tudo validado com intervalos de confian\c{c}a bootstrap de 10.000 reamostras e corre\c{c}\~ao de Holm--Bonferroni em 68 testes. V\'arias de nossas refer\^encias s\~ao preprints, refletindo um campo onde estudos fundamentais sobre o comportamento de infer\^encia de LLMs est\~ao apenas come\c{c}ando a surgir; nossa contribui\c{c}\~ao fundamenta essas observa\c{c}\~oes com o rigor estat\'istico que a pesquisa confi\'avel em IA demanda.

Acreditamos que este trabalho se alinha fortemente com o engajamento cont\'inuo da \textit{Nature Machine Intelligence} com a confiabilidade de LLMs. Kumar et al.\ examinaram recentemente quando LLMs podem ser confi\'aveis como avaliadores (\textit{Nat.\ Mach.\ Intell.}\ \textbf{8}, 173--185, 2026), e Ciriello chamou aten\c{c}\~ao para o preocupante aumento da suspei\c{c}\~ao sobre IA generativa em publica\c{c}\~oes acad\^emicas (\textit{Nat.\ Mach.\ Intell.}, 2026, \href{https://doi.org/10.1038/s42256-026-01178-z}{doi:10.1038/s42256-026-01178-z}). Seus editoriais t\^em pedido transpar\^encia em frameworks de LLM (2024) e sistemas multiagente (2026). Nossos achados falam diretamente a essas preocupa\c{c}\~oes: se as sa\'idas dos modelos mais amplamente utilizados em pesquisa n\~ao podem ser reproduzidas, a credibilidade de todo estudo derivado est\'a em risco.

Ap\'os a aceita\c{c}\~ao, depositaremos o protocolo completo e todos os 4.104 registros experimentais no Protocol Exchange e disponibilizaremos o reposit\'orio GitHub (\url{https://github.com/Roverlucas/genai-reproducibility-protocol}) sob licen\c{c}as abertas (MIT para c\'odigo, CC-BY 4.0 para dados). Este manuscrito n\~ao foi publicado em outro local e n\~ao est\'a sob considera\c{c}\~ao de nenhum outro peri\'odico. Todos os autores aprovaram a submiss\~ao.

\closing{Atenciosamente,}

\vspace{-1cm}
\noindent\textit{Em nome de todos os autores:}\\
Lucas Rover, Eduardo Tadeu Bacalhau, Anibal Tavares de Azevedo e Yara de Souza Tadano

\end{letter}
\end{document}
